\section{Introduction}

\begin{figure}[t]
    \centering
    \includegraphics[width=\textwidth]{sections/figures/teaser.pdf}
    \caption{Overview of \ourmethod{}. The target model executes tasks with a current skill, an additional frontier optimizer model converts trajectories into bounded add/delete/replace skill edits, and a held-out gate accepts only edits that improve validation performance. Accepted edits are exported as a reusable skill artifact, while rejected edits become negative feedback for later updates.}
    \label{fig:skillopt_teaser}
\end{figure}

Frontier language models are increasingly deployed as agents, from single-prompt callers to multi-step execution harnesses with tools, files, and verifiers \citep{yao2022react,schick2023toolformer,wang2023voyager,yang2024sweagent}. In such settings, domain adaptation is no longer only about model weights or prompts: it also requires improving the \emph{procedures} by which the agent gathers evidence, calls tools, follows domain conventions, and formats outputs \citep{yang2023large,khattab2023dspy}. Agent skills provide a natural interface for this procedural adaptation \citep{li2026skillsbench,jiang2026sok}: a skill is a portable natural-language artifact that packages procedures, domain heuristics, tool policies, output constraints, and failure modes, letting a frozen agent adapt through external text.

If the recurring object of adaptation is the agent's procedure, the skill document itself should be trainable. Yet weight adaptation is often unavailable for closed frontier models and expensive for open ones, while manually written or one-shot skills are brittle under a target domain or harness. Recent systems convert execution experience into reusable textual artifacts---distilling trajectory lessons, refining skill folders via failure analysis, building domain-specific skill libraries, or optimizing prompts from trajectory feedback \citep{ni2026trace2skill,alzubi2026evoskill,luo2026skillforge,shen2026skillfoundry,agrawal2025gepa}---but leave open a more basic question: if skills are the adaptation layer, how should they be optimized? Our key idea is to treat skill editing as a controllable domain-adaptation process, with the skill document as the external state, an additional frontier model as the optimizer, and training-style controls over evidence, step size, validation, and update direction.

We introduce \ourmethod{}, a text-space optimizer for agent skills. Given a target domain, an initial skill, and the model being adapted, \ourmethod{} repeatedly samples trajectory batches, analyzes successes and failures, and asks a frontier optimizer model to propose structured add/delete/replace edits. It then aggregates and ranks candidate edits under a textual learning-rate budget, applies a bounded update to the skill document, and evaluates the candidate skill on a held-out selection split before accepting it. Rejected edits are retained as negative feedback, while the epoch-wise slow/meta update preserves longer-horizon regularities. Figure~\ref{fig:skillopt_teaser} gives a schematic view of this loop. The deployed output is a compact \mbox{\texttt{best\_skill.md}} file of roughly $300$--$2{,}000$ tokens, with the adapted model and execution harness remaining fixed.

The deep-learning analogy is operational rather than decorative. Rollout and reflection batch sizes control the noise in the evidence used for each edit; the textual learning rate and schedule control how far one skill version is allowed to move from the previous one; the held-out gate plays the role of validation; and the epoch-wise slow/meta update acts like a momentum term, carrying stable editing directions across epochs. This stability is crucial: if consecutive skill revisions move too far or in inconsistent directions, rejected edits and previous accepted edits no longer provide a meaningful optimization history. With bounded, validation-gated updates, each revision remains close enough to the last one that later optimizer calls can learn from what helped, what failed, and what should be preserved.


We conduct, to our knowledge, the first systematic study of skill optimization as a domain-adaptation training method for frontier agents. We evaluate \ourmethod{} on six benchmarks covering QA, spreadsheets, documents, math, and embodied decision making, across seven target models from frontier-scale GPT to small-scale Qwen, and under three execution modes (direct chat, Codex harness, Claude Code harness). Out of 52 evaluated (model, benchmark, harness) cells, \ourmethod{} is the best or tied-best measured method on all 52. With GPT--5.5 in direct chat, it lifts SearchQA from 77.7 to 87.3, SpreadsheetBench from 41.8 to 80.7, OfficeQA from 33.1 to 72.1, DocVQA from 78.8 to 91.2, LiveMathematicianBench from 37.6 to 66.9, and ALFWorld from 83.6 to 95.5 (a $+23.5$ point average gain over no skill), and it also beats the strongest \emph{per-cell} baseline drawn from human-written, one-shot LLM, Trace2Skill, TextGrad, GEPA, and EvoSkill skills by $+5.4$ points on average. The same optimization interface is effective inside Codex-style and Claude Code-style execution loops, lifting GPT--5.5 by $+24.8$ and $+19.1$ points over no skill respectively, and outperforming EvoSkill by $+14.0$ and $+3.2$ points.

The learned artifacts also transfer beyond the exact training setting. A SpreadsheetBench skill trained on GPT--5.4 improves every smaller GPT variant we test; a Codex-trained spreadsheet skill transfers to Claude Code with a $+59.7$ point gain; and an OlympiadBench skill yields positive gains on Omni-MATH~\citep{gao2024omnimathuniversalolympiadlevel}. These transfer results are important for the paper's application value: a skill can be optimized once, audited as text, and reused across related models, harnesses, or tasks without changing model weights. Our ablations explain why this works. Bounded textual learning outperforms uncontrolled rewriting, held-out gating prevents harmful proposals from accumulating, the rejected-step buffer converts failed edits into negative feedback, and the epoch-wise slow/meta update improves long-horizon refinement without bloating the deployed skill. Finally, per-benchmark case studies show that the learned skills remain compact ($300$--$2{,}000$ tokens after only $1$--$4$ accepted edits), inspectable, and procedural rather than instance-specific.

Our contributions are as follows:
\begin{itemize}
    \item We formulate agent-skill learning as optimization over an external natural-language state and introduce \ourmethod{}, a harness-agnostic optimizer with rollout batches, reflection minibatches, add/delete/replace edits, textual learning rates, schedules, held-out acceptance, rejected-edit buffers, and epoch-wise slow/meta update.
    \item We provide a broad empirical study across six benchmarks, seven target models, and three execution harnesses, showing that \ourmethod{} is best or tied-best on 52 of 52 cells and outperforms no-skill, human-skill, one-shot LLM-skill, prompt-optimization (TextGrad, GEPA), and skill-evolution (Trace2Skill, EvoSkill) baselines under every model.
    \item We validate the optimization design through component ablations and three forms of transfer (cross-model, cross-harness, cross-benchmark), showing that the exported skill artifact is compact, reusable, and deployable without model-weight updates.
\end{itemize}
